\chapter{La Margen y el Distrito 7 (1984--2084)}

\section{Geografía urbana de La Margen}

En el año 2084, NeoBilbao no es todavía el abismo vertical de los Zócalos Industriales de 2102, pero la segregación espacial es implacable. La ciudad se organiza en tres anillos concéntricos:

\begin{enumerate}
  \item \lead{La Cúpula Dorada (Nivel Superior):} Complejo bioclimático de titanio y cristal reforzado donde residen los accionistas, ingenieros jefe y ejecutivos de Orekak-Corp. Disfrutan de aire filtrado, gravedad regulada y luz solar simulada.
  \item \lead{La Margen (Distrito 7, Nivel Medio):} La densa colmena urbana donde malvive la clase media-baja. Ingenieros junior, administrativos, sanitarios y técnicos de mantenimiento hacinados en rascacielos modulares prefabricados de hormigón poroso.
  \item \lead{El Subsuelo y la Vieja Ría (Nivel Inferior):} Dársenas industriales abandonadas, túneles de drenaje pluvial y conductos de ventilación masivos donde se refugian los proscritos, criminales y los Descartados de los experimentos corporativos.
\end{enumerate}

\section{Lugares clave del Distrito 7}

\subsection{1. La Tasca «Zulo» de Jon Ander}
Ubicada bajo el viaducto de la Línea 4 del tren mag-lev. El estruendo de los convoyes magnéticos cada seis minutos amortigua las conversaciones privadas. Jon Ander sirve cerveza sintética barata y kalimotxo fosforescente. Tras el mostrador esconde un escáner de frecuencias casero y un rifle de clavos neumáticos. Es el punto de inicio de la aventura y el refugio natural de la cuadrilla.

\subsection{2. La Estación Mag-Lev Central (D7)}
El nudo de transporte que conecta La Margen con las fábricas periféricas. Un hervidero de miles de trabajadores con abrigos de vinilo y paraguas transparentes. Cuenta con tornos de acceso biométrico que escanean la \textbf{Barik Corporativa} y calculan en tiempo real el Índice Azul de cada ciudadano. Quienes caen por debajo de 60 puntos son apartados discretamente a las cabinas de interrogatorio de los Secretas.

\subsection{3. El Dispensario de Gurutzeta Periférico}
Un hospital masificado de doce plantas donde el Celador realiza turnos interminables. Los pasillos huelen a ozono, desinfectante barato y vómito inducido por la comida sintética. En la planta baja se encuentra la unidad de psiquiatría de agudos, donde la corporación interna a ciudadanos que han sufrido «episodios de desorientación cívica».

\subsection{4. El Gaztetxe Digital «Kutxa»}
Un antiguo taller de bobinado eléctrico okupado por la Biohacker y su red de resistencia punk. Las paredes están forradas con láminas de plomo y cartón de huevos para bloquear barridos de radio. En su interior zumban decenas de ordenadores artesanales con pantallas CRT de fósforo verde, cables puenteados a la red de alta tensión y un soldador siempre humeante.

\subsection{5. Las Bocas del Colector y la Vieja Ría}
Canales de desagüe de tres metros de diámetro por donde fluye el lodo químico de la Cúpula hacia el estuario exterior. Aunque están acordonadas con rejas de hierro oxidado y carteles de peligro biológico, constituyen la única ruta practicable para descender al subsuelo sin activar los sensores de superficie.

\section{La vida bajo el Índice Azul}

El Índice Azul no es una mera calificación académica; es la moneda de cambio de la dignidad en 2084:

\begin{itemize}
  \item \lead{Créditos de oxígeno y agua:} El suministro de aire en los micro-apartamentos está tasado. Si tu Índice Azul baja por retrasos laborales o comentarios sediciosos, el caudal del respirador doméstico se reduce al mínimo de supervivencia.
  \item \lead{La cultura de la denuncia vecinal:} Los ciudadanos compiten ferozmente entre sí. Denunciar a un vecino que oculta tecnología no registrada o que consume tabaco de contrabando otorga bonificaciones inmediatas en el test del Programa Azul.
  \item \lead{El sirimiri ácido:} La condensación química procedente de los reactores de la Cúpula precipita constantemente en forma de una llovizna ácida corrosiva. Transitar por la calle sin un paraguas de polímero homologado provoca dermatitis abrasivas y arruina la ropa sintética en cuestión de semanas.
\end{itemize}
